% 2026-11-10 / v22
\section{全てのArtin切断のsmoothnessから完備基底のsmoothnessは従わない}
\label{as:sec:obstruction}

本節では、有限Artin基底での生成子の持上げと、
完備局所基底での生成性の間にある障害を明示する。
出発点となる$\mathbb P^1$上の生成対象の変形例は、
SashaのMathOverflow回答\cite{AMSasha}にある既知の例である。
以下の四次元strict dg模型、および全Artin切断の相対smoothnessと
generic fiberの非smoothnessを併記した整理の先行研究上の初出は未確認である。

この例は、非可換cubicから構成した本稿の$\Gamma_R$が非smoothであることを示すものではない。
示すのは、底上proper、滑らかな特殊ファイバー、
全ての有限Artin切断の相対smoothnessという抽象的条件だけでは、
完備基底上の相対smoothnessを結論できないという事実である。

\subsection{生成対象がgeneric fiberで生成しなくなる族}

$R=\CC[[h]]$、$K=\CC((h))$、$X=\mathbb P^1_R$とし、
$H^0(X,\mathcal O_X(1))$の基底を$s,t$とする。
ここで$s,t$は射影直線の一次sectionであり、
MFの周期変数$t$とは無関係である。
混同を避けるため、本節でのみこの記号を用いる。
二項perfect複体
\begin{equation}\label{as:eq:P}
 P_h=\bigl[\mathcal O_X(-1)\xrightarrow{\,hs\,}\mathcal O_X\bigr]
\end{equation}
をcohomological degrees $-1,0$に置く。

特殊ファイバーでは微分が零なので
\[
 P_0\simeq\mathcal O_{\mathbb P^1}\oplus
                  \mathcal O_{\mathbb P^1}(-1)[1].
\]
これは$\operatorname{Perf}(\mathbb P^1)$の古典的生成対象である。
一方、generic fiberでは$h$は可逆なので、
$s$の零点$p$について
\[
 P_h\otimes_R K\simeq\mathcal O_p.
\]
$\mathcal O_p$から有限coneと直和因子によって作る対象のsupportは$p$に含まれるため、
$\mathcal O_p$は$\operatorname{Perf}(\mathbb P^1_K)$を生成しない。
従って$P_h$も$\operatorname{Perf}(\mathbb P^1_R)$を生成しない。
実際、生成していれば、その有限生成操作を$K$へ底変更することで
generic fiberでも生成してしまう。

\begin{lemma}[各有限切断では生成する]\label{as:lem:artin-generation}
$R_N=R/(h^{N+1})$とし、$P_N=P_h\otimes_RR_N$とする。
このとき$P_N$は$\operatorname{Perf}(\mathbb P^1_{R_N})$を生成する。
\end{lemma}
\begin{proof}
射影直線の相対Beilinson分解により、perfect圏の底変更と
perfect対象間のHomの導来底変更は$R_N\to\CC$と両立する。
従って補題\ref{am:lem:generator}のInd圏における議論を、そのまま
$\operatorname{Perf}(\mathbb P^1_{R_N})$に適用できる。
具体的には、$P_N$と右直交するInd対象$M$について、
$P_N$のcompact性と底変更から
\[
 \operatorname{RHom}(P_0,M\otimes_{R_N}^{\mathbf L}\CC)=0
\]
となる。$P_0$は生成対象なので$M\otimes_{R_N}^{\mathbf L}\CC=0$である。
$R_N$の有限$h$-filtrationにより
$R_N\in\operatorname{thick}_{D(R_N)}(\CC)$だから、$M=0$となる。
これで$P_N$のcompact生成性、従ってperfect圏の古典的生成性を得る。

ここで用いた射影直線の有限模型は、対角の完全列
\[
 0\longrightarrow\mathcal O(-1,-1)
 \longrightarrow\mathcal O
 \longrightarrow\mathcal O_\Delta\longrightarrow0
\]
から得られる。従ってこの底変更の主張は、非正則なArtin基底でも成立する。
\end{proof}

\subsection{四次元strict dg代数}

\begin{proposition}[底上properな明示的dg模型]\label{as:prop:dg-model}
$\Gamma_R=\operatorname{RHom}_X(P_h,P_h)$は、次のstrict unital dg代数で表される：
\begin{equation}\label{as:eq:dg-algebra}
 \Gamma_R^0=Re_-\oplus Re_+,\qquad
 \Gamma_R^1=Rs\oplus Rt,\qquad
 \Gamma_R^q=0\quad(q\ne0,1),
\end{equation}
\[
 d(e_-)=hs,\qquad d(e_+)=-hs,\qquad d(s)=d(t)=0.
\]
$e_-,e_+$は互いに直交する冪等元で、単位は$e_-+e_+$である。
$s,t$は$e_-$から$e_+$へのdegree一の射であり、
\[
 e_+s=s=se_-,\qquad e_+t=t=te_-,
\]
その他の冪等元と射の積、および二本の射同士の積は零である。
特に$\Gamma_R$は$R$上perfectである。
この模型の形成は各$R_N$と$K$への導来底変更と両立する。
\end{proposition}
\begin{proof}
$P_h$の二つの項の間のsheaf Homは
$\mathcal O_X(-1),\mathcal O_X,\mathcal O_X(1)$の直和である。
これらの一次以上のcohomologyは零なので、
大域Homの二項複体が$\operatorname{RHom}$を実際に計算する。
degree $-1$の大域Homは$H^0(\mathcal O_X(-1))=0$であり、
degree零は二つの項のscalar endomorphism、degree一は
$H^0(\mathcal O_X(1))=Rs\oplus Rt$である。
複体の微分を$D=hs$とすると、endomorphismの微分は
$f\mapsto Df-(-1)^{|f|}fD$なので、表示した微分を得る。
積はendomorphismの合成そのものである。

この計算は有限自由$R$-加群とその積だけからなり、
同じline bundle cohomologyの消滅は各底変更後も成立する。
従って通常のtensorがこの模型の導来底変更を計算し、
底変更後の$P_h$のderived endomorphismを与える。
\end{proof}

\begin{lemma}[generic fiberの外積代数模型]\label{as:lem:generic}
$E=K[\epsilon]/(\epsilon^2)$に$|\epsilon|=1$、零微分を与える。
このとき
\[
 E\longrightarrow\Gamma_R\otimes_R K,\qquad
 1\longmapsto e_-+e_+,\qquad \epsilon\longmapsto t
\]
はdg代数のquasi-isomorphismである。
\end{lemma}
\begin{proof}
単位と積と微分を保つことは式\eqref{as:eq:dg-algebra}から直ちに分かる。
$h$が可逆なので
\[
 H^0(\Gamma_R\otimes_R K)=K(e_-+e_+),\qquad
 H^1(\Gamma_R\otimes_R K)=(Ks\oplus Kt)/Ks=K\cdot[t]
\]
であり、他のcohomologyは零である。
従って上の写像は全てのcohomologyで同型となる。
\end{proof}

\subsection{無限derived tensorによる非smoothness}

\begin{lemma}[degree一の外積代数はsmoothでない]
\label{as:lem:exterior-not-smooth}
前補題の$E$は$K$上properだが、homologically smoothではない。
\end{lemma}
\begin{proof}
augmentation $E\to K$によって$K$を左右の$E$-加群とする。
左$E$-加群のsemifree resolution
\[
 Q=\bigoplus_{n\ge0}Eu_n,\qquad |u_n|=0,\qquad
 d(u_0)=0,\quad d(u_n)=\epsilon u_{n-1}\quad(n\ge1)
\]
を考える。$Q\to K$は$u_0\mapsto1$、$u_n\mapsto0$（$n\ge1$）で与える。
$n$に関する増大filtrationはsemifree filtrationであり、
$d^2=0$は$\epsilon^2=0$による。
degree零のcycleは$Ku_0$だけであり、degree一の全ての元
$\epsilon u_n$は$d(u_{n+1})$である。
従って$Q\to K$はquasi-isomorphismである。

右からaugmentation加群$K$をtensorすると
\[
 K\otimes_E^{\mathbf L}K
 \simeq K\otimes_E Q
 =\bigoplus_{n\ge0}Ku_n
\]
となり、右辺の微分は零で、全ての$u_n$はdegree零にある。
従って$H^0(K\otimes_E^{\mathbf L}K)$は無限次元である。
もし$K$が$E$上perfectなら、このderived tensorは
有限個の$K$のshiftから有限coneと直和因子で作られ、
有限次元の有界cohomologyを持つはずなので矛盾する。

さらに、もし$E$の対角が$E^{\rm op}\otimes_KE$上perfectなら、
その有限cell表示の一方の$E$-作用へ$-\otimes_E^{\mathbf L}K$を施すと、
$K$の有限$E$-cell表示を得る。
これは$K$が$K$上perfectであることを用いる。
従ってsmoothnessは$K\in\operatorname{Perf}(E)$を含意するが、
直前でこれを否定した。
よって$E$はhomologically smoothでない。
一方、$E$自体は二次元の有界$K$-複体なのでproperである。
\end{proof}

\begin{theorem}[全Artin切断ではsmoothだが完備基底ではsmoothでない族]
\label{as:thm:obstruction}
命題\ref{as:prop:dg-model}の$\Gamma_R$は、次の全ての性質を持つ。
\begin{enumerate}
 \item $\Gamma_R$は$R$上perfectなstrict dg代数である。
 \item $\Gamma_R\otimes_R^{\mathbf L}\CC$は$\CC$上smoothである。
 \item 各$N\ge0$に対し
 $\Gamma_R\otimes_R^{\mathbf L}R_N$は$R_N$上smoothである。
 \item $\Gamma_R\otimes_R K$は$K$上smoothでなく、
 $\Gamma_R$も$R$上smoothでない。
\end{enumerate}
\end{theorem}
\begin{proof}
第一項は有限自由模型から従う。
補題\ref{as:lem:artin-generation}により$P_N$は
$\operatorname{Perf}(\mathbb P^1_{R_N})$の生成対象であり、
命題\ref{as:prop:dg-model}の底変更から
\[
 \operatorname{Perf}(\Gamma_R\otimes_R^{\mathbf L}R_N)
 \simeq\operatorname{Perf}(\mathbb P^1_{R_N})
\]
となる。
相対射影直線の対角は前述の二項線束分解を持つので、
このperfect圏は$R_N$上smoothである。
smoothnessはdg Morita同値で不変だから、第二・第三項を得る。

補題\ref{as:lem:generic}と
\ref{as:lem:exterior-not-smooth}からgeneric fiberはsmoothでない。
$\Gamma_R$が$R$上smoothなら、そのperfect対角の表示を
flatな底変更$R\to K$へ移してgeneric fiberもsmoothとなるため、矛盾する。
従って第四項も成立する。
\end{proof}

\begin{remark}[この反例の適用範囲]
$R_N\in\operatorname{thick}_{D(R_N)}(\CC)$はArtin基底での生成子の持上げを支える。
しかし$R=\CC[[h]]$は$D(R)$で$\CC$の有限coneと直和因子から作れない。
実際、そのような対象を$K$へtensorすると零になるが、$R\otimes_RK=K\ne0$である。
上の族では、この違いが生成性とsmoothnessの実際の破綻として現れる。

従って、全Artin比較の両立性だけから、完備基底上の生成性や相対smoothnessを補ってはならない。
可換側がsmooth properな射影族であることだけでも、変形した対象の生成性は保証されない。
中心三次非可換商の残余圏へこの原理を適用するには、
当該windowと生成対象についての有限生成証明、
または対角の具体的なperfect表示を別途用いる。
本例は、その特別な主張を否定してはいない。

なお本例の$R$上のcohomologyは
$H^0(\Gamma_R)=R$、$H^1(\Gamma_R)=Rt\oplus(R/hR)s$であり、
cohomologyの底上flatnessは満たさない。
この例が否定する仮定の範囲を、証明していない別の仮定へ広げてはならない。
例\ref{cg:ex:artin-singularity}はMFと絶対特異圏の区別を示すものであり、
本節のsmoothnessの持上げ障害とは別の反例である。
\end{remark}
